\documentclass[../../main.tex]{subfiles}
\begin{document}

\section{范畴与函子}\label{section:ch03-范畴与函子}

范畴论为切空间和微分提供了另一种有用的视角。本节中我们概述范畴论的基本定义。我们在本书中并不怎么使用这个理论，但提到它，是因为它为谈论我们将遇到的许多数学结构提供了一种方便而有力的语言。

一个范畴 $\mathcal{C}$ 由下列事物组成：
\begin{itemize}
\item 一个类 $\mathrm{Ob}(\mathcal{C})$，其元素称为 $\mathcal{C}$ 的对象；
\item 一个类 $\mathrm{Hom}(\mathcal{C})$，其元素称为 $\mathcal{C}$ 的态射；
\item 对每个态射 $f\in\mathrm{Hom}(\mathcal{C})$，两个对象 $X,Y\in\mathrm{Ob}(\mathcal{C})$，分别称为 $f$ 的源和靶；
\item 对每个三元组 $X,Y,Z\in\mathrm{Ob}(\mathcal{C})$，一个称为复合的映射
\begin{align*}
\mathrm{Hom}_{\mathcal{C}}(X,Y)\times\mathrm{Hom}_{\mathcal{C}}(Y,Z)\to\mathrm{Hom}_{\mathcal{C}}(X,Z),
\end{align*}
记作 $(f,g)\mapsto g\circ f$，其中 $\mathrm{Hom}_{\mathcal{C}}(X,Y)$ 表示所有源为 $X$、靶为 $Y$ 的态射所成类。
\end{itemize}

这些态射需要满足下列公理：
\begin{enumerate}[(i)]
\item \textbf{结合性}：$(f\circ g)\circ h=f\circ(g\circ h)$。
\item \textbf{恒等态射的存在性}：对每个对象 $X\in\mathrm{Ob}(\mathcal{C})$，存在一个恒等态射 $\mathrm{Id}_{X}\in\mathrm{Hom}_{\mathcal{C}}(X,X)$，使得对所有 $f\in\mathrm{Hom}_{\mathcal{C}}(Y,X)$ 有 $\mathrm{Id}_{Y}\circ f=f=f\circ\mathrm{Id}_{X}$。
\end{enumerate}

一个态射 $f\in\mathrm{Hom}_{\mathcal{C}}(X,Y)$ 称为 $\mathcal{C}$ 中的一个同构，若存在一个态射 $g\in\mathrm{Hom}_{\mathcal{C}}(Y,X)$ 使 $f\circ g=\mathrm{Id}_{Y}$ 且 $g\circ f=\mathrm{Id}_{X}$。

\begin{example}[范畴]\label{example:ch03-范畴}
在人们"在自然界中"遇到的大多数范畴里，对象是带某种额外结构的集合，态射是保持该结构的映射，而复合律和恒等态射是显而易见的。本书中（隐含或显式）出现的一些这种类型的范畴列在下面。在每种情形，我们通过给出其对象和其态射来描述该范畴。
\begin{itemize}
\item $\mathsf{Set}$：集合与映射；
\item $\mathsf{Top}$：拓扑空间与连续映射；
\item $\mathsf{Man}$：拓扑流形与连续映射；
\item $\mathsf{Man}^{b}$：带边拓扑流形与连续映射；
\item $\mathsf{Diff}$：光滑流形与光滑映射；
\item $\mathsf{Diff}^{b}$：带边光滑流形与光滑映射；
\item $\mathsf{Vec}_{\mathbb{R}}$：实向量空间与实线性映射；
\item $\mathsf{Vec}_{\mathbb{C}}$：复向量空间与复线性映射；
\item $\mathsf{Grp}$：群与群同态；
\item $\mathsf{Ab}$：交换群与群同态；
\item $\mathsf{Rng}$：环与环同态；
\item $\mathsf{CRng}$：交换环与环同态。
\end{itemize}
\end{example}

还有一些重要范畴，其对象除（可能）其他结构外，是带指定基点的集合。一个带点集是一个有序对 $(X,p)$，其中 $X$ 是一个集合，$p$ 是 $X$ 的一个元素。其他带点对象，如带点拓扑空间或带点光滑流形，类似地定义。若 $(X,p)$ 和 $(X',p')$ 是带点集（或拓扑空间等），一个映射 $F:X\to X'$ 称为一个带点映射，若 $F(p)=p'$；此时我们写 $F:(X,p)\to(X',p')$。下面是带点对象范畴的一些重要例子。
\begin{itemize}
\item $\mathsf{Set}_{*}$：带点集与带点映射；
\item $\mathsf{Top}_{*}$：带点拓扑空间与带点连续映射；
\item $\mathsf{Man}_{*}$：带点拓扑流形与带点连续映射；
\item $\mathsf{Diff}_{*}$：带点光滑流形与带点光滑映射。
\end{itemize}

我们对范畴中的对象类和态射类使用"类"一词而非"集合"，因为在某些范畴中它们"太大"，不能被视作集合。例如，在范畴 $\mathsf{Set}$ 中，$\mathrm{Ob}(\mathsf{Set})$ 是所有集合所成类；任何把它本身当作一个集合来处理的尝试，都会导致众所周知的集合论 Russell 悖论。（更多见 [LeeTM, 附录~A] 或几乎任何一本集合论书。）即使对象类和态射类不构成集合，我们仍使用诸如 $X\in\mathrm{Ob}(\mathcal{C})$ 和 $f\in\mathrm{Hom}(\mathcal{C})$ 的记号来表示 $X$ 是 $\mathcal{C}$ 中的一个对象、$f$ 是 $\mathcal{C}$ 中的一个态射。一个其中 $\mathrm{Ob}(\mathcal{C})$ 和 $\mathrm{Hom}(\mathcal{C})$ 都是集合的范畴称为小范畴，而其中每个态射类 $\mathrm{Hom}_{\mathcal{C}}(X,Y)$ 都是集合的范畴称为局部小范畴。上面列出的所有范畴都是局部小范畴但不是小范畴。

若 $\mathcal{C}$ 和 $\mathcal{D}$ 是范畴，则一个从 $\mathcal{C}$ 到 $\mathcal{D}$ 的协变函子是这样一条规则 $\mathcal{F}$：它给每个对象 $X\in\mathrm{Ob}(\mathcal{C})$ 指派一个对象 $\mathcal{F}(X)\in\mathrm{Ob}(\mathcal{D})$，并给每个态射 $f\in\mathrm{Hom}_{\mathcal{C}}(X,Y)$ 指派一个态射 $\mathcal{F}(f)\in\mathrm{Hom}_{\mathcal{D}}(\mathcal{F}(X),\mathcal{F}(Y))$，使得恒等和复合被保持：
\begin{align*}
\mathcal{F}(\mathrm{Id}_{X})=\mathrm{Id}_{\mathcal{F}(X)};\qquad \mathcal{F}(g\circ h)=\mathcal{F}(g)\circ\mathcal{F}(h).
\end{align*}

我们还需要考虑反转态射的函子：一个从 $\mathcal{C}$ 到 $\mathcal{D}$ 的反变函子是这样一条规则 $\mathcal{F}$：它给每个对象 $X\in\mathrm{Ob}(\mathcal{C})$ 指派一个对象 $\mathcal{F}(X)\in\mathrm{Ob}(\mathcal{D})$，并给每个态射 $f\in\mathrm{Hom}_{\mathcal{C}}(X,Y)$ 指派一个态射 $\mathcal{F}(f)\in\mathrm{Hom}_{\mathcal{D}}(\mathcal{F}(Y),\mathcal{F}(X))$，使得
\begin{align*}
\mathcal{F}(\mathrm{Id}_{X})=\mathrm{Id}_{\mathcal{F}(X)};\qquad \mathcal{F}(g\circ h)=\mathcal{F}(h)\circ\mathcal{F}(g).
\end{align*}

\begin{exercise}\label{exercise:ch03-3.27}
证明：任意（协变或反变）从 $\mathcal{C}$ 到 $\mathcal{D}$ 的函子把 $\mathcal{C}$ 中的同构映为 $\mathcal{D}$ 中的同构。
\end{exercise}

协变函子的一个平凡例子是从任意范畴到其自身的恒等函子：它把每个对象和每个态射映到自身。另一个例子是遗忘函子：若 $\mathcal{C}$ 是一个其对象为带某种额外结构的集合、其态射为保持该结构的映射的范畴（如例~\cref{example:ch03-范畴} 第一部分所列、除 $\mathsf{Set}$ 本身外的所有范畴），则遗忘函子 $\mathcal{F}:\mathcal{C}\to\mathsf{Set}$ 给每个对象指派其底集，给每个态射指派那个被视作集合之间映射的同一映射。

当我们把"不变量"与数学对象的类联系起来时，会出现更有趣的函子。例如，基本群是一个从 $\mathsf{Top}_{*}$ 到 $\mathsf{Grp}$ 的协变函子。习题~2-10 的结果表明，存在一个从 $\mathsf{Diff}$ 到 $\mathsf{Vec}_{\mathbb{R}}$ 的反变函子，其定义为：给每个光滑流形 $M$ 指派向量空间 $C^{\infty}(M)$，给每个光滑映射 $F:M\to N$ 指派由 $F^{*}(f)=f\circ F$ 定义的线性映射 $F^{*}:C^{\infty}(N)\to C^{\infty}(M)$。

本章的讨论给了我们其他一些重要的函子例子。首先，切空间函子是一个从带点光滑流形范畴 $\mathsf{Diff}_{*}$ 到实向量空间范畴 $\mathsf{Vec}_{\mathbb{R}}$ 的协变函子。它给每个带点光滑流形 $(M,p)$ 指派切空间 $T_{p}M$；给每个带点光滑映射 $F:(M,p)\to(N,F(p))$ 指派微分 $\mathrm{d}F_{p}$。这是一个函子这一事实，正是命题~\cref{proposition:ch03-微分性质}(b)、(c) 的内容。类似地，我们可以把指派 $M\mapsto TM$、$F\mapsto\mathrm{d}F$（把每个光滑流形送到其切丛、把每个光滑映射送到其整体微分）看作一个从 $\mathsf{Diff}$ 到其自身的协变函子，称为切函子。

\end{document}